\chapter{イントロダクション}

この本は理系全般，特に数学に興味のあるの大学生・社会人に向けて
情報学とプログラミングの入門となることをまとめた．
特に情報学とプログラミングの歴史に関係したことと
そこに表れる数学に関連したこととから説き起こし，
すでに知っていると思われる高校の情報Iのレベルより
更に高いレベルへの入門を目指した．

このような本をまとめた背景を述べたい．
学生に接しているとCLI=Command Line Interfaceを触ったことがなく，
開発環境をローカルに構築したことがない，VSCodeすら触ったことの無い
学生が多数いることがわかる．
またコロナ禍以降，スマートフォン中毒者もよく見かけるようになった．
スマートフォンは時間を意味なく消費し，
目や耳や首や肩などを始めとして体の不調を引き起こし，
脳の発達を邪魔し，アルゴリズムが世界観をも壊している．
更に，これからは生成AIに依存・代替される人間が増えることが予想されている．

これらに対応するため，これからの学生・社会人には
\begin{itemize}
   \item
      スマートフォンとインターネットと生成AIに適切な距離を置く
   \item
      CLIを使い(\textcolor{red}{生成AIは既に使いこなしている}ので，
      逆に生成AIを使いこなすためには必須である)こなす
   \item
      必要に応じて自力で開発環境を構築し，ツールがつかえるようになる
   \item
      幾つかのプログラミング言語でプログラミングができる
   \item
      生成AIにMarkdown などを上手く使って適切に構造化した指示が出来る
   \item
      生成AIにより拡散された偽情報を見抜ける
\end{itemize}
ことが大切になると思われる．

そこでこの本では
CLIの演習を行い，
デジタルデータの仕組みを述べ，取り扱いを説明する．
次に論理回路やチューリングマシンなど計算論の基礎への説明を与え，
インターネットやWEBの仕組みやセキュリティーについて説明する．
最後にはデジタルデトックスについても論じている．
付録には具体的で身近なプログラミング言語
(例えばExcel/VBA, \LaTeX, markdown, JavaScript, Shell Script, git)
のプログラミング演習となるような資料を幾つか上げたので，
成るべく多くのものにチャレンジし，技術を深掘りしてほしい．
また，演習問題も多数あげたので自分の理解の確認に使って貰えると有り難い．
授業のレポートで友達のものやネットに上がっているものを単にコピペしても
理解も技術も深まらないように，生成AIの答えを貼り付けても理解も技術も深まりません．
それらを参考にしても構わないが，大凡のアイデアを理解したら，
自分なりにプログラミングの本質を理解し，アイデアを検討しなおし，
適当な関数にまとめるとか，クラスにまとめるとかして下さい．
コードが読みやすいようにソースを適切にまとめ，分かり易い変数の名前を選ぶ
とか，適切なコメントを入れるとかしてください．
そうすれば理解も技術も深めて自分のものにすることができる．
%生成AIについてもできる限り取り上げている．
全体を合わせると，
上に上げた項目に対処できる力に一歩でも近づけるものが増えると期待している．

なおこの本では MacOSの英語表示を使って説明をしている．
Windows でも大体は同じことが成り立つが，
慣れないと対応が分かりづらい部分もあると思う．
ネット等で調べたり生成AIに聞いて各自対応してほしい．

この本の図のソースファイルや関連した動画やアプリを
\href{https://snow00two.github.io/mySite/}{私のサイト}に用意したので
それらを動かしながら理解を深めてもらいたい．
